\begin{korollar}{Existenz unabhängiger (\textsl N(0,1),
\textsl n)--Tupel}{ExistenzNormalNTupel}
Seien $n \in \N$ und $\alpha \in \N^n$.\\
Dann existiert ein unabhängiges $(N(0,1),n)$--Tupel zu $(n, \alpha)$.
\end{korollar}
\begin{proof}
Sei $N\coloneqq  \Sum{i=1}n \alpha_i$.\\
Nach~\refclaim{ExistenzFVerteilterZufallsgrössen}
{ExistenzFVerteilterZufallsgrössen2} gibt es einen Wahrscheinlichkeitsraum
$(\Omega', \mfa', \Prim')$ und $\eta: \Omega' \rightarrow \R^N$ unabhängig und
$N(0,1)$--verteilt.\\
Definiere für alle $i \in \haken n$:
\[s_i \coloneqq  \left(\sum_{k=1}^{i-1}\alpha_i\right)+1\text{, }
  t_i \coloneqq  \sum_{k=1}^{i}\alpha_i \text{ und }
 \gamma_i\coloneqq \left(\eta_j \right)_{j = s_i}^{t_i}\text.\]
Seien nun $k, l \in \haken n$ mit $k \neq l$. Dann sind $\gamma_k$ und
$\gamma_l$ vollständig unabhängig als Teilfamilien einer unabhängigen Familie
und damit ist \[\bigg((\Omega',\mfa',\Prim'),(\gamma_i)_{i=1}^n, (\alpha_i)_{i=1}^n\bigg)\]
ein $(N(0,1),n)$--Tupel zu $(n, \alpha)$.
\end{proof}